\section*{Chapter \ref{ch:dv:autocallables} --- Autocallables}

\begin{solution}{exo:dv:autocallables:1}
A zero-coupon bond repaying par (long); a European down-and-in put struck at 100\% with barrier 60\% (short); a strip of \omterm{def:dv:barriers-and-digitals:digital}{digital
options} paying the coupon when the index is at or above 70\% on each date, with memory (long); the early-redemption feature, a
strip of digitals at 100\% that end the note (short, since the issuer chooses nothing but the investor loses the remaining
coupons and the upside above par).
\end{solution}

\begin{solution}{exo:dv:autocallables:2}
The investor finances the coupon mainly by selling the knock-in put. Under the smile that put is priced at the high volatilities of
low strikes, so it is worth more and pays for more coupon; the smile's lower short-dated at-the-money volatilities also change the
call probabilities. Flat volatility undervalues the put and so the coupon it can finance: 3.14\% against 6.18\% a year.
\end{solution}

\begin{solution}{exo:dv:autocallables:3}
$(98-95.99)/2.60=0.78\%$ a quarter, 3.1\% a year.
\end{solution}

\begin{solution}{exo:dv:autocallables:4}
The issuer is long the knock-in put. Higher dividends lower the forward, which makes the put more valuable: the issuer gains, so it is
long dividends (the chapter's example: a 1\% yield moves the note's value by $-0.20$ for the investor, $+0.20$ for the issuer). It sells
the exposure in dividend futures or swaps.
\end{solution}

\begin{solution}{exo:dv:autocallables:5}
The note's delta to the investor grows as the underlying approaches the barrier, because the knock-in put gains value quickly there;
the issuer, short the note, holds that delta in the underlying and buys as it falls (and sells as it rises). At the knock-in the note
becomes a long position with a delta of about one, far below the delta just above the barrier, and the issuer sells the difference
into the fall.
\end{solution}

\begin{solution}{exo:dv:autocallables:6}
Just above the barrier the note's value falls from about par plus the coupon to the knocked-in value over a distance that shrinks like
$\sigma\sqrt\tau$: the less time is left, the steeper the drop and the larger the delta before the knock-in. After it the delta is about
one whatever the time left. The difference, the cliff, grows like $1/\sqrt\tau$: 97, 254 and 497 million for a year, three months and one
month in the chapter's example.
\end{solution}

\begin{solution}{exo:dv:autocallables:7}
The fair coupon rises from 1.55\% to 1.91\% a quarter (7.62\% a year), and the knock-in probability from 6.3\% to 10.4\%: a barrier monitored
every day is breached by paths that recover before maturity, so the investor sells a more valuable put and is paid more for it.
\end{solution}

\begin{solution}{exo:dv:autocallables:8}
The hedges are local: at a daily knock-in the delta jumps and the desk must sell at once, into a falling market, alongside every other
issuer with notes at the same level; a gap through the level leaves the sale at a worse price (chapter 15). The \omterm{def:dv:greeks-and-the-hedging-pnl:greeks}{vega} and skew of the book
change abruptly too. The event needs a gap reserve, limits on concentration at barrier levels, and pre-positioning.
\end{solution}

\begin{solution}{pb:dv:autocallables:1}
\textbf{1.} 100.44 (standard error 0.06); knock-in probability 15.6\%.
\textbf{2.} A daily down-and-in put struck at 100\% with barrier 75\%, and the upside, called away at 103\%; they finance the accruing
coupon.
\textbf{3.} Knock-out: par plus 15\% a year to the knock-out date. Neither: par plus 30\% at maturity. Knock-in without a later knock-out:
par times the final performance, capped at par.
\textbf{4.} It can be breached by a path that recovers before maturity, and it is monitored every day, including the days of sharp falls
and gaps.
\textbf{5.} The \omterm{def:dv:autocallables:snowball}{snowball}, 15\% against 6.18\% a year: its knock-in is higher (75\% against 60\%) and monitored daily, so the investor sells
a much more valuable put.
\textbf{6.} 7.63.
\textbf{7.} 1.00.
\textbf{8.} $(7.63-1.00)\times0.75\times100$ million $=497$ million.
\textbf{9.} 254 million and 97 million.
\textbf{10.} Bought the underlying, more and more as the delta grew, from about one at 100\% to 7.6 near the barrier.
\textbf{11.} One barrier ten times as large: a sale of the order of five billion concentrated at one level.
\textbf{12.} About 0.6 times notional with twelve months left and about 6 times with one month.
\textbf{13.} They are the cheapest and fastest hedge for an index, and the flow is executed there first; basis and activity show it.
\textbf{14.} Limit issuance on the same underlying in view of its trading volume and the maximum delta the hedge can reach.
\textbf{15.} Outstanding notional by underlying, the distribution of knock-in levels and of maturities, and the issuers' dealer gamma.
\textbf{16.} A reserve for the gap and the market impact of the sale: size the sale, estimate its cost from depth and impact models
(Book 8), and charge it, weighted by the probability of reaching the barrier near maturity.
\textbf{17.} Reduce the delta ahead of the barrier (under-hedge the last points), buy put spreads around the barrier, or pre-sell part of the
knock-in delta in a controlled way.
\textbf{18.} The investor, whose note now loses with the index; the issuer's P\&L depends on how well the sale was executed.
\textbf{19.} 497 million of the underlying (at the barrier level) per 100 million of notes: the delta falls from 7.63 to 1.00.
\textbf{20.} Many issuers hedge notes with the same barrier and maturities, so the hedge sale at the level is the whole market's.
\end{solution}

\begin{solution}{iq:dv:autocallables:1}
A bond plus: short a knock-in put (the capital at risk), long coupon digitals (with memory in a Phoenix), short the upside through the
early call. The coupon is financed by the put premium and by the call cutting the note short in good scenarios.
\iqlookfor{the decomposition and who is short what.}
\end{solution}

\begin{solution}{iq:dv:autocallables:2}
Start with \omterm{def:dv:local-volatility:model}{local volatility}, which fits the vanillas that price the terminal distributions; check with a local-stochastic model for the
forward-skew-dependent coupons and calls; reserve for the difference. Flat volatility underprices the knock-in put and mispricing the
digitals: it roughly halves the fair coupon in the chapter's example.
\iqlookfor{\omterm{def:dv:local-volatility:model}{local volatility}, its limits, and the direction of the flat-volatility error.}
\end{solution}

\begin{solution}{iq:dv:autocallables:3}
Long \omterm{def:dv:greeks-and-the-hedging-pnl:greeks}{vega} (short end and maturity), long skew, long dividends, and for worst-ofs short correlation. The issuer sells volatility (options,
\omterm{def:dv:variance-swaps-and-volatility-derivatives:vs}{variance swaps}), skew (risk reversals) and dividends (futures, swaps) to the market, and holds correlation or offsets it with index against
single-stock options.
\iqlookfor{the four exposures and the recycling instruments.}
\end{solution}

\begin{solution}{iq:dv:autocallables:4}
Use common random numbers for the bumped prices; many paths; smooth the digital features (replace the digital by a narrow call spread, as an
overhedge would) or use pathwise and likelihood-ratio estimators (chapter 23); bucket the \omterm{def:dv:greeks-and-the-hedging-pnl:greeks}{Greeks} rather than bump each node; check stability
across seeds.
\iqlookfor{common random numbers and a treatment of the discontinuities.}
\end{solution}

\begin{solution}{iq:dv:autocallables:5}
Aggregate outstanding notional by underlying, knock-in level and maturity; for each level compute the change in the issuers' delta when it is
crossed (the cliff) as a function of time to maturity; compare with the underlying's and futures' daily volume and depth.
\iqlookfor{notional times barrier and maturity concentration times dealer gamma.}
\end{solution}

\begin{solution}{iq:dv:autocallables:6}
Estimate the market's sale at the level; reduce my own delta ahead of it and buy protection (puts or put spreads below the level); widen quotes
for new notes; prepare to provide liquidity after the level breaks if the sale overshoots; tell risk management and size the gap reserve.
\iqlookfor{anticipating the flow, protection, and liquidity after the break.}
\end{solution}
